\section{Related Work}
\label{sec:related}

\paragraph{Variational optimization and optimal control.}
Continuous-time optimization has been studied through gradient flows,
accelerated ODEs, Bregman Lagrangians, mirror descent, and natural-gradient
geometry
\cite{su2016differential,wibisono2016variational,krichene2015accelerated,
beck2003mirror,amari1998natural}.  Our controlled variable is the metric
itself, the terminal set is determined by preconditioning quality, and the
value measures the intrinsic deformation required to reach that set.

Min-plus semigroups, dynamic programming, and Hamilton--Jacobi equations are
classical in optimal control \cite{fleming2006controlled}.  Geometric control
also supplies the classical second variation, Jacobi fields, and conjugate
point analysis \cite{agrachev2004control}.  We use these ingredients as the
analytic foundation for a different reduction target: an entire metric path
ending at a preconditioning-quality set.  The new response statements are the
intervention-cube Green representation and its bordered hard-target analogue.

\paragraph{Parametric optimization and Jacobi response.}
Reduced Hessians, variable projection, solution-map sensitivity, and implicit
differentiation eliminate optimized variables through Schur complements
\cite{golub1973differentiation,bonnans2000perturbation,
dontchev2014implicit,amos2017optnet,blondel2022efficient,
franceschi2018bilevel}.  We do not claim the generic finite-dimensional Schur
formula as new.  The contribution is its coercive path-Hilbert realization,
the derivation of that coercivity from Hadamard geometry on an entire
intervention neighborhood, the joint bulk/terminal Green operator, and the
resulting finite- and arbitrary-order interaction formulas.  For an active hard
terminal inequality, the closest general framework is parametric constrained
optimization and generalized equations
\cite{bonnans2000perturbation,dontchev2014implicit}; our bordered theorem
identifies the derivative explicitly as a Jacobi--KKT boundary operator and
specializes it to the moving condition-number projection.

Jacobi fields and index forms classically describe variations of geodesics.
Here the Jacobi operator contains kinetic, curvature, bulk-potential, and
terminal Robin terms.  Nonpositive curvature contributes a nonnegative term
to the index form, allowing global control even when the convex potentials
have linear negative growth.

The finite-dimensional companion theory already contains the generic hidden
response, Schur reduced Hessian, affine-intervention negative Gram operator,
and its pair finite-difference integral
\cite{li2026optimizationgeometrodynamics}.  The local formulas in
\cref{thm:path-schur,cor:path-interaction} are their Hilbert-path lift and are
used here as infrastructure.  The additions proved in this paper are the
Hadamard hypotheses that derive global path existence, uniqueness, and
cube-wide coercivity; the strong Green--Jacobi boundary problem that unifies
bulk and Robin forcing; arbitrary prescribed finite-order path recursion; and
the bordered Jacobi--KKT response of a moving hard target.  These conclusions
are absent from the finite-dimensional companion reduction theorem.

\paragraph{Positive-definite geometry and structured preconditioning.}
The affine-invariant geometry of positive-definite matrices supplies the
distance, exponential map, and nonpositive curvature used in the SPD
specialization
\cite{bhatia2007positive,pennec2006riemannian,higham2008functions}.
Geodesic convex analysis and metric projection on Hadamard spaces are
developed in
\cite{sra2015conic,lim2004best,bacak2014convex}.
Smoothness on a fixed simple-eigenvalue stratum and loss of classical
differentiability at spectral multiplicities follow the spectral-function
theory in \cite{lewis1996derivatives,lewis2001twice,lewis2003eigenvalue}.

Classical equilibration and condition-number minimization optimize an endpoint
scaling or matrix approximation
\cite{vandersluis1969equilibration,marechal2009optimizing,
lu2011minimizing,diamond2016stochastic,tanaka2013positive,
qu2024optimal,dogan2025geodesicpreconditioning}.  These works clarify the
endpoint benchmark addressed by \cref{thm:full-spd-benchmark}.  RDGC adds the
initial metric and an admissible path geometry; its hard-target theorem
differentiates the projection path in addition to the endpoint condition
number.

Diagonal, block, and Kronecker preconditioners arise in adaptive and
second-order optimization
\cite{duchi2011adagrad,kingma2014adam,shazeer2018adafactor,
martens2015kfac,grosse2016kfc,gupta2018shampoo,anil2020scalable}.
These algorithms motivate restricted metric families, but a concrete
optimizer-to-path statement must additionally specify the optimizer state,
damping, bias correction, step map, and interpolation.  RDGC is conditional
on that realization and measures the geometric effort once it has been
declared.

\paragraph{Finite interaction decompositions.}
Boolean and factorial interactions can be represented by finite differences,
M\"obius coefficients, and higher-order attribution indices
\cite{pukelsheim2006optimal,dasgupta2015causal,dhamdhere2020shapley,
janizek2021explaining}.  The result here is a path-geometric representation:
the pair coefficient is a cube-face integral of an inverse-Jacobi Gram kernel,
conditional effects use the same kernel on subfaces, and all higher
coefficients follow from one Green recursion.


